% ============================================================
% §4 The GAINS Algorithm and Theoretical Analysis
% §4.1 Offline scheduling
% §4.2 Online control
% §4.3 SDE foundations (derives location-scale from Taylor expansion)
% §4.4 Scheduling optimality (water-filling + Jensen gap, now as corollaries)
% ============================================================
\section{The GAINS Algorithm}
\label{sec:algorithm}

GAINS schedules noise trajectory search at two granularities:
offline coarse allocation (\cref{sec:offline}) sets per-timestep budgets from
a small calibration run, and online control (\cref{sec:online}) adapts those
budgets per prompt.
The theoretical analysis proceeds in two stages.
\cref{sec:sde-foundations,sec:scheduling-optimality} treat the
\emph{special case} where the local operator is pure random search
(i.i.d.\ noise draws), deriving a gain factorization
$G_t(K)=\sigma_t\,a_K$ and its scheduling consequences
(water-filling, Jensen gap).
\cref{sec:general-local-search} then \emph{generalizes} these results
to an abstract class of local search operators satisfying four regularity
conditions, covering the $\epsilon$-greedy and local perturbation methods
used in our experiments.

% -----------------------------------------------------------
\subsection{Offline Coarse Time Scheduling}
\label{sec:offline}

We consider the common setting where the sampler maintains \emph{a single latent and a single noise}
$(x_t, \varepsilon_t)$ at each timestep $t$, but performs \emph{multiple local search
iterations} before committing to $x_{t-1}$.
Under this restriction, the global scheduling problem reduces to a sequence of
\emph{per-timestep decisions}: at each $t$, given the history of search so far
and the remaining budget, should we continue local search at this timestep,
or stop and move on to $t-1$?

An inherent difficulty arises: when deciding at timestep $t$, we do not know:
(1)~the potential improvement that additional search at $t$ could still bring,
and (2)~the potential improvement that could be obtained by saving NFE and
spending it on later timesteps.
The scheduling strategy below addresses both aspects from complementary
angles.

For each model and dataset, we first run \emph{offline} profiling on a
small set of prompts or images.
We execute noise search with a simple local operator $\mathcal{L}_t$ and a uniform,
relatively large iteration count at every timestep, logging verifier
scores after each iteration.
From these logs we compute, for each timestep $t$, the
average score gain from local search and its variability across examples.
Aggregating over many runs reveals which segments of the
denoising trajectory respond most to noise refinement and
are therefore most ``valuable'' to search.
Offline profiling reveals a clear model-dependent pattern:
for Stable Diffusion (SD), gains concentrate on early timesteps,
whereas for EDM they cluster in a contiguous middle segment.

We then partition timesteps into a small number of groups,
e.g., a \emph{high-value} region and a \emph{low-value} region, assigning a
larger share of the total NFE budget to the high-value region.
Within each region, we derive a coarse per-step budget $\{K_t\}$ (more
iterations in high-value steps, fewer in low-value steps).
The resulting schedule, tailored to the model and dataset but fixed across
prompts, captures \emph{where search matters most}.

% -----------------------------------------------------------
\subsection{Online Fine-Grained Scheduling}
\label{sec:online}

At test time, an \emph{online} controller refines how the
pre-assigned budget $\{K_t\}$ is spent on each prompt.
At timestep $t$, local search runs up to $K_t$ iterations, but the
controller can stop early or, within a small window, modulate how
aggressively to spend the budget.

The controller monitors two statistics from local search
at timestep $t$:
(i)~the \emph{score gain} of local search, and
(ii)~the \emph{score variance} across candidates within this step.
Let $s_t^{(j)}$ be the best verifier score at step $t$ after $j$ local search
iterations, and let $\{s_{t,\text{cand}}^{(j,i)}\}_i$ denote the scores of all
candidate noises evaluated in iteration $j$.
We define the incremental gain
\begin{equation}
  g_t^{(j)} = s_t^{(j)} - s_t^{(j-1)}
  \label{eq:gain-increment}
\end{equation}
and the per-iteration variance
\begin{equation}
  \operatorname{Var}_t^{(j)} = \operatorname{Var}_i\bigl(s_{t,\text{cand}}^{(j,i)}\bigr).
  \label{eq:var-iteration}
\end{equation}
Here $g_t^{(j)}$ measures whether search still finds
better noises at step $t$, while $\operatorname{Var}_t^{(j)}$ reflects the
verifier's sensitivity to noise perturbations at this timestep.

The online policy follows a simple rule:
if recent gains $g_t^{(j)}$ remain positive and score variance
$\operatorname{Var}_t^{(j)}$ is large, the timestep is both important
and unsaturated, so local search continues to budget $K_t$ (or a small upper margin).
Conversely, when gains diminish (recent $g_t^{(j)}$ near
zero or negative) and candidate variance is small, local
search at this timestep is exhausted and further NFE is better reserved
for subsequent timesteps; the controller stops early and advances to $t-1$.
As a further safeguard, whenever an iteration produces a negative gain
($g_t^{(j)} < 0$), the controller reverts the noise pivot to the
previous best candidate rather than accepting a worse solution.
This \emph{revert-on-negative} rule ensures that local search at any
timestep is monotonically non-decreasing in verifier score.
All decisions rely solely on statistics observed up to step
$t$, with no access to future timesteps.
To match realized compute exactly to the prescribed total NFE budget,
we reconcile over-/under-spending from adaptive stopping:
extra NFE consumed in the high-value region is offset by allocating
additional iterations to the first few low-value steps, while any NFE deficit
is corrected by trimming iterations from the last low-value steps.

\paragraph{Overall procedure.}
A sampling run draws initial noise $x_T \sim \mathcal{N}(0, I)$ and
reads the coarse per-step budget $\{K_t\}$ from the offline schedule.
For $t = T, \dots, 1$, the local search operator
$\mathcal{L}_t$ explores around $x_t$ while the online controller
monitors gains $\{g_t^{(j)}\}$ and variances
$\{\operatorname{Var}_t^{(j)}\}$ to decide whether to continue or stop,
subject to budget $K_t$.
The chosen noise at step $t$ feeds into the base sampler
$F_\theta$ to produce $x_{t-1}$, and the process repeats until $x_0$.
Under a fixed NFE budget, this two-level scheduling concentrates
computation on timesteps that are both globally and instance-wise valuable,
avoiding wasted search where noise perturbations barely affect
the verifier.
\cref{alg:offline_online_budget} gives the complete procedure.

\begin{algorithm}[t]
\caption{GAINS: Global Adaptive Inference-time Noise Scheduling}
\label{alg:offline_online_budget}
\begin{algorithmic}[1]
\STATE \textbf{Input:} timesteps $\{t_T,\dots,t_1\}$, total budget $B$, offline allocation $\{K_t\}$ ($\sum_t K_t=B$),
slack $\delta$, gain window $W_g$, threshold coefs $(\beta_g,\beta_\sigma)$, black-box $\textsc{LocalIter}(t)$.
\STATE \textbf{Output:} realized per-step iterations $\{\widehat K_t\}$.

\STATE $R \leftarrow B$; \quad $\mathcal{H}_g\leftarrow\emptyset$ (historical gains); \quad $\mathcal{H}_\sigma\leftarrow\emptyset$ (historical variances).
\FOR{$t=T,T\!-\!1,\dots,1$}
    \STATE $K \leftarrow \min(K_t, R)$; \quad $K_{\max}\leftarrow \min(K+\delta, R)$; \quad $j_{\text{watch}}\leftarrow \max(1, K-\delta)$.
    \STATE $s^{(0)} \leftarrow -\infty$; \quad $\mathcal{H}\leftarrow\emptyset$ (recent gains); \quad $\widehat K_t \leftarrow 0$.
    \FOR{$j=1,2,\dots,K_{\max}$}
        \STATE $(s^{(j)}, \sigma^{(j)}) \leftarrow \textsc{LocalIter}(t)$
        \IF{$s^{(j)} < s^{(j-1)}$}
            \STATE $s^{(j)} \leftarrow s^{(j-1)}$ \hfill \COMMENT{revert-on-negative: keep previous best}
        \ENDIF
        \STATE $g^{(j)} \leftarrow s^{(j)} - s^{(j-1)}$; push $g^{(j)}$ into $\mathcal{H}$ and keep last $W_g$.
        \STATE $\widetilde g^{(j)} \leftarrow \mathrm{Mean}(\mathcal{H})$.
        \STATE $\widehat K_t \leftarrow \widehat K_t + 1$; \quad $R \leftarrow R - 1$.
        \IF{$R=0$} \STATE \textbf{break} \ENDIF
        \IF{$j \ge j_{\text{watch}}$ \AND $|\mathcal{H}_g|>0$ \AND $|\mathcal{H}_\sigma|>0$}
            \STATE $g_{\mathrm{th}}\leftarrow \beta_g \cdot \mathrm{Mean}(\mathcal{H}_g)$;\quad
                   $\sigma_{\mathrm{th}}\leftarrow \beta_\sigma \cdot \mathrm{Mean}(\mathcal{H}_\sigma)$.
            \IF{$\widetilde g^{(j)} < g_{\mathrm{th}}$ \AND $\sigma^{(j)} < \sigma_{\mathrm{th}}$}
                \STATE \textbf{break}
            \ENDIF
        \ENDIF
    \ENDFOR
    \STATE append $\mathrm{Mean}(\mathcal{H})$ to $\mathcal{H}_g$; append $\sigma^{(\widehat K_t)}$ to $\mathcal{H}_\sigma$.
\ENDFOR
\STATE \textbf{return} $\{\widehat K_t\}$.
\end{algorithmic}
\end{algorithm}

% -----------------------------------------------------------
\subsection{SDE Foundations for the Score Distribution}
\label{sec:sde-foundations}

Both stages of GAINS rely on the per-step score
distribution $\mathcal{D}_t$ exhibiting location-scale structure.
We first derive this structure for the special case of
\emph{pure random search} (i.i.d.\ noise draws), where the
analysis reduces to Gaussian order statistics.
The extension to general operators appears in
\cref{sec:general-local-search}.

The reverse-time SDE in Euler form with step size $h$ and diffusion
coefficient $g_t$ reads:
\begin{equation}
  X_{t-h} \;=\; X_t + b_t(X_t, c)\,h + g_t\sqrt{h}\,\xi,
  \qquad \xi \sim \mathcal{N}(0, I),
  \label{eq:euler-sde}
\end{equation}
where $b_t(X_t, c)$ denotes the drift (determined by the learned score network)
and $g_t\sqrt{h}\,\xi$ the stochastic perturbation that noise search
operates on.
The \emph{verifier score function}
$\Psi_t(x; c) := v(D_\theta(x, t{-}h); c)$
maps a latent $x$ to the verifier score of its denoised prediction,
so the score from one noise candidate equals
\begin{equation}
  S_t \;:=\; \Psi_t(X_{t-h};\, c).
  \label{eq:score-from-sde}
\end{equation}
Let $\bar{X}_{t-h} := X_t + b_t(X_t, c)\,h$ denote the deterministic Euler
point, the latent that would result without noise injection.

\begin{assumption}[Verifier smoothness]
\label{asm:smooth}
The composed verifier score function
$\Psi_t(\cdot\,; c) \in C^2(\mathbb{R}^d)$, and there exists a constant
$L_t < \infty$ such that
$\|\nabla^2 \Psi_t(x; c)\|_{\mathrm{op}} \le L_t$
for all $x$ in a neighborhood of the line segment between $\bar{X}_{t-h}$
and $X_{t-h}$, almost surely.
\end{assumption}

\begin{proposition}[Score Taylor expansion]
\label{prop:taylor}
Under \cref{asm:smooth}, conditional on $(X_t, c)$,
\begin{equation}
  S_t \;=\; \Psi_t(\bar{X}_{t-h};\, c)
        \;+\; g_t\sqrt{h}\,
              \bigl\langle \nabla\Psi_t(\bar{X}_{t-h};\, c),\, \xi \bigr\rangle
        \;+\; R_t,
  \label{eq:taylor-score}
\end{equation}
where the remainder satisfies
$|R_t| \le \tfrac{1}{2}\,L_t\,g_t^2\,h\,\|\xi\|^2$ almost surely,
and hence
$\mathbb{E}[\,|R_t|\mid X_t, c\,] \le \tfrac{1}{2}\,L_t\,g_t^2\,h\,d$.
\end{proposition}

\begin{proof}
By construction,
$X_{t-h} = \bar{X}_{t-h} + g_t\sqrt{h}\,\xi$.
A second-order Taylor expansion of $\Psi_t(\cdot\,; c)$ around
$\bar{X}_{t-h}$ gives
\[
  \Psi_t(X_{t-h}; c)
  = \Psi_t(\bar{X}_{t-h}; c)
    + \bigl\langle \nabla\Psi_t(\bar{X}_{t-h}; c),\,
        X_{t-h} - \bar{X}_{t-h} \bigr\rangle
    + R_t,
\]
where
$R_t = \tfrac{1}{2}(X_{t-h} - \bar{X}_{t-h})^\top
  \nabla^2\Psi_t(\widetilde{X}; c)\,
  (X_{t-h} - \bar{X}_{t-h})$
for some $\widetilde{X}$ on the segment joining $\bar{X}_{t-h}$ and $X_{t-h}$.
Substituting $X_{t-h} - \bar{X}_{t-h} = g_t\sqrt{h}\,\xi$ yields
\eqref{eq:taylor-score}.
The remainder bound follows from
$|R_t| \le \tfrac{1}{2}\|\nabla^2\Psi_t(\widetilde{X}; c)\|_{\mathrm{op}}
  \,\|g_t\sqrt{h}\,\xi\|^2
  \le \tfrac{1}{2}L_t\,g_t^2\,h\,\|\xi\|^2$.
Taking conditional expectations and using
$\mathbb{E}[\|\xi\|^2] = d$ completes the proof.
\end{proof}

\cref{prop:taylor} decomposes each candidate score into a deterministic
baseline $\Psi_t(\bar{X}_{t-h}; c)$, a Gaussian linear term with variance
governed by $g_t\sqrt{h}\,\|\nabla\Psi_t\|$, and a higher-order remainder
that vanishes as $h$ shrinks.
This decomposition yields a location-scale characterization.

\begin{theorem}[Location-scale structure of per-step scores]
\label{thm:loc-scale}
Define
\begin{equation}
  \mu_t \;:=\; \Psi_t(\bar{X}_{t-h};\, c),
  \qquad
  \sigma_t \;:=\; g_t\sqrt{h}\,\|\nabla\Psi_t(\bar{X}_{t-h};\, c)\|.
  \label{eq:mu-sigma-sde}
\end{equation}
Then, conditional on $(X_t, c)$:
\begin{enumerate}
\item[\textbf{(i)}]
  There exists $Z_t \sim \mathcal{N}(0,1)$ such that
  $S_t = \mu_t + \sigma_t\,Z_t + R_t$.
\item[\textbf{(ii)}]
  As $g_t\sqrt{h} \to 0$,
  \begin{equation}
    \operatorname{Var}(S_t \mid X_t, c)
    \;=\; \sigma_t^2 + O_d(g_t^3\,h^{3/2}).
    \label{eq:var-approx}
  \end{equation}
\item[\textbf{(iii)}]
  Let $S_t^{(1)}, \ldots, S_t^{(K)}$ be $K$ i.i.d.\ copies of $S_t$
  generated by i.i.d.\ noise draws $\xi_1, \ldots, \xi_K$.
  Then
  \begin{equation}
    G_t(K)
    \;:=\; \mathbb{E}\Bigl[\max_{1 \le k \le K} S_t^{(k)}
           \,\Big|\, X_t, c\Bigr] - \mu_t
    \;=\; \sigma_t\,a_K + O_{K,d,L_t}(g_t^2\,h),
    \label{eq:gain-sde}
  \end{equation}
  where $a_K := \mathbb{E}[\max(Z_1, \ldots, Z_K)]$ with
  $Z_k \stackrel{\mathrm{iid}}{\sim} \mathcal{N}(0,1)$.
  In particular,
  \begin{equation}
    G_t(K) \;=\; g_t\sqrt{h}\,\|\nabla\Psi_t(\bar{X}_{t-h};\, c)\|\,a_K
              \;+\; O_{K,d,L_t}(g_t^2\,h).
    \label{eq:gain-explicit}
  \end{equation}
\end{enumerate}
\end{theorem}

\begin{proof}
For part~(i), when $\|\nabla\Psi_t(\bar{X}_{t-h}; c)\| > 0$, write
$\langle \nabla\Psi_t, \xi \rangle
= \|\nabla\Psi_t\|\,Z_t$
where $Z_t := \langle \nabla\Psi_t / \|\nabla\Psi_t\|, \xi \rangle$.
Since $\nabla\Psi_t / \|\nabla\Psi_t\|$ is a fixed unit vector conditional on
$(X_t, c)$ and $\xi \sim \mathcal{N}(0, I)$, the projection $Z_t$ is standard
normal: $Z_t \sim \mathcal{N}(0,1)$.
(When $\nabla\Psi_t(\bar{X}_{t-h}; c) = 0$, we have $\sigma_t = 0$ and the
linear term vanishes; the result holds trivially with $Z_t$ defined
arbitrarily.)

For part~(ii), expand
$\operatorname{Var}(S_t \mid X_t, c)
= \operatorname{Var}(\sigma_t Z_t + R_t)
= \sigma_t^2 + 2\,\operatorname{Cov}(\sigma_t Z_t, R_t) + \operatorname{Var}(R_t)$.
By Cauchy--Schwarz,
\[
  |\operatorname{Cov}(\sigma_t Z_t, R_t)|
  \le \sqrt{\operatorname{Var}(\sigma_t Z_t)}\,
      \sqrt{\operatorname{Var}(R_t)}
  \le \sigma_t\,\sqrt{\mathbb{E}[R_t^2]}.
\]
Using \cref{prop:taylor},
$|R_t| \le \tfrac{1}{2}L_t g_t^2 h \|\xi\|^2$, hence
\[
  \mathbb{E}[R_t^2 \mid X_t, c]
  \le \tfrac{1}{4}L_t^2 g_t^4 h^2\,\mathbb{E}[\|\xi\|^4]
  = \tfrac{1}{4}L_t^2 g_t^4 h^2\,d(d+2),
\]
because $\xi\sim\mathcal{N}(0,I)$ implies
$\mathbb{E}[\|\xi\|^4]=d(d+2)$.
Therefore
\[
  |\operatorname{Cov}(\sigma_t Z_t, R_t)|
  \le \tfrac{1}{2}L_t\sqrt{d(d+2)}\,\sigma_t g_t^2 h
  = O_d(g_t^3 h^{3/2}),
\]
with the implicit constant depending on the conditional local geometry
through $L_t$ and $\|\nabla\Psi_t\|$.
Likewise,
\[
  \operatorname{Var}(R_t)
  \le \mathbb{E}[R_t^2]
  \le \tfrac{1}{4}L_t^2 g_t^4 h^2\,d(d+2)
  = O_d(g_t^4 h^2),
\]
which is lower order than the covariance term as $g_t\sqrt{h}\to 0$.

For part (iii), write
$S_t^{(k)} = \mu_t + \sigma_t Z_k + R_t^{(k)}$
for each candidate $k$.
Then
$\max_k S_t^{(k)} = \mu_t + \max_k (\sigma_t Z_k + R_t^{(k)})$.
Using the elementary inequality
$|\max_k (a_k + b_k) - \max_k a_k| \le \max_k |b_k|$,
we obtain
$|\max_k S_t^{(k)} - (\mu_t + \sigma_t \max_k Z_k)|
\le \max_k |R_t^{(k)}|$.
Since $\max_k |R_t^{(k)}| \le \sum_k |R_t^{(k)}|$,
taking conditional expectations gives
$\mathbb{E}[\max_k |R_t^{(k)}| \mid X_t, c]
\le K \cdot \tfrac{1}{2} L_t g_t^2 h\,d$.
Therefore
$G_t(K) = \sigma_t\,a_K + O_{K,d,L_t}(g_t^2 h)$,
where the implicit constant in $O_{K,d,L_t}(\cdot)$ depends on $K$, the
dimension~$d$, and the Hessian bound~$L_t$ from \cref{asm:smooth}, but not
on the diffusion coefficient $g_t$ or the verifier gradient
$\nabla\Psi_t$.
Substituting $\sigma_t = g_t\sqrt{h}\,\|\nabla\Psi_t(\bar{X}_{t-h}; c)\|$
yields~\eqref{eq:gain-explicit}.
\end{proof}

\begin{remark}
\cref{thm:loc-scale} identifies the precise quantities governing timestep
sensitivity:
$\sigma_t = g_t\sqrt{h}\,\|\nabla\Psi_t(\bar{X}_{t-h}; c)\|$
multiplies the diffusion coefficient $g_t$, the step size
$\sqrt{h}$, and the verifier gradient norm at the
deterministic Euler point.
Timesteps with large $g_t$ (strong noise injection) or large
$\|\nabla\Psi_t\|$ (steep verifier landscape) benefit most from noise
search, explaining the non-uniform sensitivity patterns observed in offline profiling.
\end{remark}

\begin{corollary}[SDE-based approximate allocation]
\label{cor:sde-alloc}
Under \cref{asm:smooth}, with step size $h$ and diffusion
coefficients $\{g_t\}_{t=1}^T$, for any fixed budget allocation
$\{K_t\}_{t=1}^T$,
\begin{equation}
  \sum_{t=1}^{T} G_t(K_t)
  \;=\; \sum_{t=1}^{T} g_t\sqrt{h}\,
        \|\nabla\Psi_t(\bar{X}_{t-h}; c)\|\,a_{K_t}
  \;+\; \sum_{t=1}^{T} O_{K_t,d,L_t}(g_t^2\,h).
  \label{eq:aggregate-gain}
\end{equation}
Dropping the second-order remainder reduces the offline allocation problem
to
\begin{equation}
  \max_{\{K_t\}_{t=1}^T:\;\sum_t K_t = B}\;
  \sum_{t=1}^{T} g_t\sqrt{h}\,
  \|\nabla\Psi_t(\bar{X}_{t-h}; c)\|\,a_{K_t}.
  \label{eq:alloc-sde}
\end{equation}
\end{corollary}

Problem~\eqref{eq:alloc-sde} is a concave resource allocation (each
$a_{K_t}$ grows concavely in $K_t$) with
per-step weights $\sigma_t = g_t\sqrt{h}\,\|\nabla\Psi_t\|$.
It recovers the abstract formulation~\eqref{eq:alloc} in
\cref{sec:scheduling-optimality} below, now with an explicit,
SDE-derived expression for the per-step sensitivity $\sigma_t$.

% -----------------------------------------------------------
\subsection{Scheduling Optimality}
\label{sec:scheduling-optimality}

Under pure random search, \cref{thm:loc-scale} enables a sharp
characterization of the optimal budget allocation.
Because the gain factorizes as
$G_t(K) = \sigma_t\,a_K + O_{K,d,L_t}(g_t^2 h)$ (\cref{thm:loc-scale}(iii)),
global scheduling reduces, up to controlled remainders, to
\begin{equation}
  \max_{K_1,\dots,K_T}\;\sum_{t=1}^{T}\sigma_t\,a_{K_t}
  \quad\text{s.t.}\;\sum_{t=1}^{T}K_t=B,\;K_t\ge 1.
  \label{eq:alloc}
\end{equation}

\begin{proposition}[Offline optimality under independence]
\label{prop:offline}
Suppose $\sigma_t$ does not depend on the prompt or search history
(i.e., $\sigma_t = \sigma_t$ for all $c, x_t$).
Then:
\begin{enumerate}
\item[\textbf{(i)}]
  Problem~\eqref{eq:alloc} decomposes across prompts: any optimal
  allocation $\{K_t^*\}$ is simultaneously optimal for every instance.
\item[\textbf{(ii)}]
  The optimal allocation equalizes marginal gains:
  $\sigma_t\,a_{K_t^*}' = \lambda^*$ for all~$t$ with $K_t^* > 1$,
  where $\lambda^*$ is the shadow price of the budget constraint.
  Consequently, $K_t^*$ is increasing in~$\sigma_t$ (water-filling).
\item[\textbf{(iii)}]
  $\sum_t \sigma_t\,a_{K_t^*} > \sum_t \sigma_t\,a_{B/T}$ whenever
  $\{\sigma_t\}$ are not all equal, and the gap grows with the
  dispersion of~$\{\sigma_t\}$.
\end{enumerate}
\end{proposition}

\begin{remark}
\cref{prop:offline} justifies offline profiling.
Part~(i) guarantees that under the independence condition, a single fixed
allocation optimizes every prompt simultaneously, so profiling on a small calibration
set generalizes.
Part~(ii) reveals \emph{water-filling} structure: timesteps with larger
$\sigma_t = g_t\sqrt{h}\,\|\nabla\Psi_t\|$ (where noise perturbations
shift the verifier score most) receive proportionally more budget.
Part~(iii) quantifies the non-uniform advantage: the more
heterogeneous the per-step sensitivities, the larger the gain
over uniform scheduling.
\end{remark}

\begin{proposition}[Online adaptation under prompt dependence]
\label{prop:online}
Suppose $\sigma_t = \sigma_t(c, x_t)$ varies with the prompt and search
history.
Define the oracle value
$V^*(\boldsymbol{\sigma})=\max_{\{K_t\}:\sum_t K_t=B}
\sum_t\sigma_t\,a_{K_t}$
for variance profile $\boldsymbol{\sigma}=(\sigma_1,\dots,\sigma_T)$.
Then:
\begin{enumerate}
\item[\textbf{(i)}]
  $V^*$ is convex in $\boldsymbol{\sigma}$.
  The expected regret of any fixed allocation satisfies
  \begin{equation}
    \mathbb{E}_{\boldsymbol{\sigma}}\!\bigl[V^*(\boldsymbol{\sigma})\bigr]
    \;-\; V^*\!\bigl(\bar{\boldsymbol{\sigma}}\bigr) \;\ge\; 0,
    \label{eq:jensen-gap}
  \end{equation}
  where $\bar{\boldsymbol{\sigma}}=\mathbb{E}[\boldsymbol{\sigma}]$,
  with equality whenever the same allocation is optimal almost surely on
  the support of $\boldsymbol{\sigma}$.
\item[\textbf{(ii)}]
  For integer budgets, the one-step marginal gain obeys
  \begin{equation}
    \Delta_t(K)
    \;:=\;
    G_t(K{+}1)-G_t(K)
    \;=\;
    \underbrace{\sigma_t}_{\text{sensitivity}}
    \;\cdot\; \underbrace{\Delta a_K}_{\text{exhaustion}}
    \;+\; O_{K,d,L_t}(g_t^2 h),
    \qquad
    \Delta a_K:=a_{K+1}-a_K .
    \label{eq:marginal-decomp}
  \end{equation}
  The optimal stopping condition is $\Delta_t(K)\le\lambda^*$.
  The running mean $\bar{g}$ upper-bounds the active discrete marginals
  up to the same $O_{K,d,L_t}(g_t^2 h)$ remainder, so
  $\beta_g\bar{g}$ serves as an adaptive estimate of~$\lambda^*$.
  The gain condition $\widetilde{g}_t^{(j)}<\beta_g\bar{g}$ tests the
  leading-order product $\sigma_t\,\Delta a_K$ against the
  shadow-price scale~$\lambda^*$; the
  variance condition is an empirical proxy for low sensitivity rather
  than an exact test of $\sigma_t$.
  Requiring both conditions prevents premature stopping at
  high-value timesteps.
\end{enumerate}
\end{proposition}

\begin{remark}
\cref{prop:online} motivates the online controller.
Part~(i) shows that $V^*$, as a pointwise maximum of linear functions, is convex in the variance
profile, so by Jensen's
inequality any fixed allocation tuned to the population mean
$\bar{\boldsymbol{\sigma}}$ incurs a gap
(\cref{eq:jensen-gap}) growing with cross-instance variability
of~$\boldsymbol{\sigma}$.
Equality holds whenever the same allocation remains optimal across the
support of $\boldsymbol{\sigma}$; a constant profile is only a
sufficient special case.
Part~(ii) decomposes the marginal gain~\eqref{eq:marginal-decomp}
into two factors, each monitored by a separate
threshold in the online controller.
Low gain with high variance ($\Delta a_K$ small, $\sigma_t$ large) signals
a responsive but exhausted timestep where a lucky draw may still yield a large jump, so stopping is
premature.
Moderate gain with low variance ($\Delta a_K$ moderate, $\sigma_t$ small)
signals inherent insensitivity where gains plateau
regardless, yet search remains productive at this iteration count.
The variance statistic used by the controller should therefore be read
as a proxy for the latent sensitivity $\sigma_t$, justified by the
location-scale variance approximation in \cref{thm:loc-scale}(ii),
rather than as an exact measurement of~$\sigma_t$ itself.
Only when \emph{both} factors are small is stopping unambiguously
justified, as the two-threshold diagram in \cref{fig:stopping} shows.
\end{remark}

\begin{figure}[t]
\centering
\includegraphics[width=0.75\linewidth]{figures/fig_stopping_decision.pdf}
\caption{Idealized online stopping decision in the $(\sigma_t, a_K')$ plane.
The gain threshold (horizontal dashed) tracks exhaustion, while the
vertical threshold should be read as an empirical-variance proxy for
low sensitivity rather than an exact test of $\sigma_t$.
partition the space into four regions.
Stopping is triggered only when \emph{both} conditions are met (shaded region),
corresponding to the conjunction in \cref{prop:online}(ii).
The hyperbola $\sigma_t\cdot a_K'=\lambda^*$ shows the true optimality boundary;
the rectangular proxy is conservative by design.}
\label{fig:stopping}
\end{figure}

The proofs of \cref{prop:offline,prop:online} are given in \cref{sec:proofs}.
